\section{Operations Executive Briefing and Live Department Channels}
\label{sec:pilot-operations-executive-briefing}
\raggedright

\subsection{Purpose and delivered behaviour}

The Operations Command Centre now implements a governed executive operating layer rather than a
set of decorative department tabs.  The left pane remains the founder's direct conversation with
Operations Pilot.  The right pane contains six persistent COO-to-Department-Pilot channels:
Sales, Marketing, Finance, Support, People, and Products \& Services.  Each channel records COO
instructions, departmental updates, morning minutes, evidence provenance, model provenance, and
approval state as a durable business record.

The implementation provides the following behaviour:
\begin{itemize}[leftmargin=*]
\item a daily executive briefing scheduled for 08:00 in the configured business timezone;
\item a separate evidence-bounded director role for every participating department;
\item review of previous-day results, the current day's plan, outstanding work, risks, target
      variance, dependencies, and decisions required from the COO;
\item inclusion of published Board meeting summaries, Board mandates, delegated actions, and the
      current business pulse;
\item dated minutes written into every Department Pilot channel;
\item a briefing summary written into the main Operations Pilot conversation;
\item a persistent action register for COO follow-up;
\item staged feedback for material risks or target variance so it can be surfaced to the COO,
      CEO, and Board without silently changing Board decisions;
\item optional challenge by a second, separately accepted model; and
\item exact approval gating for any later external action.
\end{itemize}

The feature is provider-neutral.  It is not fixed to Qwen, GPT-OSS, OpenAI, or any other model.
The primary reasoner is the model selected and accepted through the Vault at the time the briefing
or channel turn runs.

\subsection{Runtime architecture}

The implementation separates cognition from authority.  A model may supply bounded departmental
judgement, but AION owns source selection, role assignment, validation, persistence, scheduling,
audit records, action state, and approval enforcement.  The effective execution path is:

\begin{enumerate}[leftmargin=*]
\item The signed desktop session fixes the actor and workspace.  Browser input cannot select a
      different business record.
\item AION identifies the requested department and loads only its authorised evidence projection.
\item AION loads the latest published Board context and the bounded recent channel history.
\item The Vault resolver selects a released model.  Non-selectable local cards and disconnected
      providers are rejected.
\item AION gives the model a department-specific role card, evidence pack, Board context, output
      schema, and explicit non-execution rules.
\item The returned JSON is parsed and validated.  Unsupported or malformed output is rejected.
\item AION writes the validated position, evidence identifiers, model receipt, and approval state
      into the durable channel.
\item Proposed actions enter an internal open-action register.  They are not treated as completed
      work and receive no external execution authority.
\end{enumerate}

In compact form, the authority boundary is:
\begin{center}
\texttt{Vault model proposal}\\
$\downarrow$\\
\texttt{AION validation and persistence}\\
$\downarrow$\\
\texttt{governed Department Pilot work}\\
$\downarrow$\\
\texttt{exact approval for external change}
\end{center}

No renderer, model response, or department channel message can mint an approval or directly invoke
an external connector.

\subsection{Department role separation}

Each Department Pilot is invoked separately with a distinct executive mandate.  This prevents a
single undifferentiated prompt from presenting the same generic advice under six different names.

\begin{longtable}{p{.18\linewidth}p{.76\linewidth}}
\textbf{Department} & \textbf{Bounded executive mandate}\\ \hline
Sales & Protect revenue quality, progress qualified pipeline, remove conversion blockers, and
state the next accountable commercial action.\\
Marketing & Protect efficient demand generation, distinguish measured attribution from
assumptions, and coordinate campaigns with Sales capacity.\\
Finance & Protect cash, margin, and financial control; distinguish ledger evidence from forecasts
and surface overdue or unreconciled exposure.\\
Support & Protect customer outcomes, service levels, and closure quality; surface complaints,
overdue replies, and cross-department blockers.\\
People & Protect team capacity, accountability, and safe access while excluding unnecessary
sensitive personal detail.\\
Products \& Services & Protect customer value, delivery feasibility, and unit economics while
separating recorded product facts from assumptions.\\
\end{longtable}

Each model response must use the supplied evidence and return a structured position containing
\texttt{summary}, \texttt{yesterday\_results}, \texttt{today\_plan},
\texttt{outstanding}, \texttt{risks}, \texttt{target\_variance},
\texttt{decisions\_needed}, and \texttt{actions}.  A department cannot change a Board target from
its channel; it can only report variance and request a decision.

\subsection{Evidence and Board context}

Department evidence is obtained from the existing provider projections rather than assembled by
the user interface.  Representative source boundaries include:
\begin{itemize}[leftmargin=*]
\item Sales: the Sales revenue spine and Boardroom snapshot;
\item Finance: the business financial model and Finance/Sales invoice ledger;
\item Marketing: Marketing department intelligence and Boardroom actions;
\item Support: the governed support-case ledger;
\item People: the organisation-authority model and authorised availability records; and
\item Products \& Services: the business operating model, offerings, inventory metrics, and unit
      economics.
\end{itemize}

The Board context is a bounded projection of the latest published meeting records, delegated
department actions, and business pulse.  Raw credentials, secrets, access tokens, private keys,
and unbounded provider payloads are excluded before model invocation.  Evidence source identifiers
are retained with the resulting message so the UI can display the number of sources and the audit
record can identify the authority used.

If evidence is absent, the retrieval state is recorded as unavailable or not published.  If the
selected model is unavailable, AION records the known source boundary and explicitly marks
executive judgement and target variance as unknown.  It does not generate a plausible but
unsupported meeting discussion.

\subsection{Primary and critic model routing}

The primary model is resolved dynamically from the workspace and Vault.  The stored public receipt
contains only bounded fields such as provider, model, source, release status, and failure reason;
Vault secrets are never copied into the conversation.

A separately accepted critic model may be configured under
\texttt{executive\_briefing\_config.critic\_model\_selection}.  The critic must resolve successfully
and must be different from the primary model.  It receives the bounded departmental positions and
Board context and may challenge unsupported claims, cross-department conflicts, or decisions that
need executive attention.  It cannot execute work.

The briefing records one of two explicit reasoning modes:
\begin{itemize}[leftmargin=*]
\item \texttt{single\_model\_role\_separated}: departments were considered through separate role
      invocations, but no independent second model participated; or
\item \texttt{two\_model\_challenge}: a distinct accepted model challenged the departmental
      positions.
\end{itemize}

This disclosure prevents one local model speaking to itself from being misrepresented as an
independent executive debate.

\subsection{Scheduling, recovery, and duplicate prevention}

The default schedule is daily at 08:00 in \texttt{Europe/Madrid}.  The local desktop backend checks
for due executive work once per minute while Tessaris is running.  The Operations page also checks
when opened and every five minutes while it remains visible.  The two paths use the same local-date
deduplication rule, so a missed or delayed check can be recovered without producing duplicate
minutes.

A normal scheduled run returns the existing briefing when that local date has already been
published.  The visible \emph{Run morning briefing} control deliberately uses a replacement run,
allowing the founder to regenerate that day's minutes after evidence or model availability changes.
The replacement does not grant external execution authority.

The desktop application must be running for its local scheduler to operate.  An optional always-on
worker node may perform the same protected check when the Mac application is not running; the
briefing service itself remains idempotent by business-local date.

\subsection{Conversation and audit persistence}

Executive messages are stored in
\texttt{operational\_runtime\_summary.shared\_conversations}.  The new record format is
\texttt{workspace\_conversation\_v2}.  Each message includes:
\begin{itemize}[leftmargin=*]
\item message identifier, sequence, department, sender, role, and message type;
\item UTC creation time and bounded structured metadata;
\item evidence source identifiers and public model selection metadata;
\item \texttt{approval\_gated} and \texttt{external\_writes\_performed};
\item the previous record hash; and
\item a canonical hash of the current record.
\end{itemize}

Relevant message types include \texttt{coo\_instruction},
\texttt{department\_update}, \texttt{morning\_minutes}, and
\texttt{executive\_briefing\_summary}.  The newest 1,200 workspace turns are retained in the active
record.  Clients request a bounded recent page rather than loading the entire history into the
renderer.

Each completed meeting is also stored in \texttt{executive\_briefings} with its local date,
meeting identifier, department positions, Board context, critic receipt, action register, Board
feedback, reasoning mode, safety state, and canonical briefing hash.  The active store retains the
newest 120 meeting records.

\subsection{Actions, follow-up, and escalation}

Every proposed departmental action becomes an entry in
\texttt{executive\_action\_register} with an open state, originating meeting, department, creation
time, and \texttt{external\_execution\_allowed=false}.  This gives Operations Pilot a durable basis
for later questions such as whether Sales reached its target or whether Support obtained a customer
reply.

Material target variance and departmental risks are copied into
\texttt{executive\_board\_feedback} with state
\texttt{staged\_for\_coo\_ceo\_board}.  Staging makes the issue available to the executive layer; it
does not rewrite approved Board minutes, change a target, or claim Board acceptance.

The COO may send a follow-up from any department channel.  The protected turn includes that
department's role card, current evidence, Board context, and the eight most recent bounded channel
messages.  The COO instruction and Department Pilot response are then appended to the same durable
stream.

\subsection{Protected API surface}

All executive endpoints are under the signed AION Flow desktop session:
\begin{longtable}{p{.30\linewidth}p{.64\linewidth}}
\textbf{Endpoint suffix} & \textbf{Purpose}\\ \hline
\texttt{executive-status} & Returns the local schedule, whether today's briefing is due, the latest
briefing receipt, and bounded channel summaries.\\
\texttt{morning-briefing} & Runs or deduplicates the six-department meeting and returns its signed
workspace result.\\
\texttt{channel-history} & Returns the bounded persistent stream for exactly one supported
department.\\
\texttt{channel-turn} & Records a COO instruction and creates an evidence-bounded reply from the
selected Department Pilot.\\
\end{longtable}

The possession-bound session supplies the person and workspace.  Payloads cannot select another
workspace.  Unsupported department identifiers and empty messages are rejected with validation
errors.

\subsection{Desktop interface}

The Operations Command Centre now displays:
\begin{itemize}[leftmargin=*]
\item the live schedule and timezone;
\item whether the meeting used role separation or a two-model challenge;
\item a manual morning-briefing control;
\item six Department Pilot tabs;
\item a persistent, scrollable message stream per department;
\item sender and timestamp for each entry;
\item message type, evidence-source count, public model identity, and approval state; and
\item a COO composer addressed to the selected Department Pilot.
\end{itemize}

The previous ``Mission bridge not connected'' placeholder has been removed.  Selecting a tab now
loads its protected business record, and sending a message invokes the governed department turn
rather than changing local screen text only.

\subsection{Failure handling and safety guarantees}

The following rules are enforced by the service rather than by model instruction alone:
\begin{itemize}[leftmargin=*]
\item A missing or unavailable model produces an explicit evidence-only record; it does not trigger
      canned business advice.
\item Invalid model JSON is rejected and replaced with an honest unavailable-model result.
\item Missing evidence remains unknown and cannot be converted into a target result.
\item A model response cannot claim an external write, issue an approval hash, or select an
      unregistered department.
\item Morning briefings are read-and-prepare operations.  They cannot send email, update CRM,
      publish advertising, move money, contact employees, or mutate connected services.
\item External work must be compiled into a registered Department Pilot capability and pass through
      AION's server-side exact approval gateway.
\item A critic model is optional and cannot silently replace the user's selected primary model.
\end{itemize}

\subsection{Implementation locations}

The principal implementation files are:
\begin{itemize}[leftmargin=*]
\item \texttt{backend/modules/aion\_business/runtime/operations\_executive\_briefing\_service.py}:
      role cards, model routing, evidence packaging, meeting execution, persistence, action staging,
      and Board feedback;
\item \texttt{backend/api/workflow\_capsule\_router.py}: signed executive status, briefing,
      history, and turn endpoints;
\item \texttt{backend/desktop\_app.py}: once-per-minute due-work scan while Tessaris is running;
\item \texttt{desktop/mac/src/app.js}: live executive-channel interface, protected request handling,
      automatic recovery checks, and COO composer; and
\item \texttt{docs/OPERATIONS\_EXECUTIVE\_BRIEFING.md}: concise operational and technical reference.
\end{itemize}

The installed application is \texttt{/Applications/Tessaris.app}.  The pre-install recovery copy is
\texttt{/Applications/Tessaris.app.pre-executive-briefing-v1}.

\subsection{Verification}

Focused automated tests cover six-department minute publication, local-date deduplication,
COO-instruction and Department-Pilot response persistence, Products \& Services identifier
normalisation, the distinct critic-model path, visible desktop controls, protected route wiring,
removal of the obsolete placeholder, and background scheduler integration.  Seven focused tests
passed.  Python compilation, Electron JavaScript syntax, scoped whitespace checks, and the installed
application's macOS signature verification also passed.

The installed UI was exercised directly.  Dated Sales and Finance minutes remained visible after an
application restart.  During that verification the selected OpenAI model adapter was unavailable;
the system correctly retained the authorised evidence identifiers and marked judgement and target
variance as unknown instead of fabricating an executive discussion.  This is the intended degraded
mode and does not change the provider-neutral design.
